\chapter{Feeling Worthless}
\index{Feeling Worthless}

\section*{Definition and Overview}
Feeling worthless is a painful, persistent sense of having little value or importance, often accompanied by harsh self-criticism, guilt, and the belief that one is a burden to others. Occasional low moments are part of normal human experience, but when feelings of worthlessness are intense, frequent, or long-lasting, they are a significant warning sign. Worthlessness is a core symptom of depression and is closely linked to anxiety, trauma, and other mental health conditions. It is also strongly associated with suicidal thinking, so it must always be taken seriously. With the right help, these feelings can lift, and this chapter explains the causes, what to do, and where to find support.

\section*{Epidemiology}
Feelings of worthlessness are extremely common and affect people of all ages, backgrounds, and life circumstances. They are one of the most frequently reported symptoms in depression, which affects around 1 in 6 people at some point in their lives, and they are more common in people facing social isolation, financial stress, chronic illness, and relationship difficulties. Young people report high rates of self-critical and worthless feelings, and these are important risk markers for self-harm and suicide. Because people often hide these feelings, the true number affected is higher than official figures suggest.

\section*{Aetiology and Pathophysiology}
\begin{itemize}
    \item \textbf{Depression:} worthlessness is a hallmark symptom, driven by changes in brain chemistry, negative thought patterns, and reduced motivation
    \item \textbf{Anxiety disorders:} chronic anxiety erodes self-confidence and fuels self-criticism
    \item \textbf{Trauma and abuse:} past emotional, physical, or sexual abuse shapes beliefs that the self is flawed or to blame
    \item \textbf{Life stressors:} relationship breakdown, job loss, financial difficulty, grief, and caring stress can trigger or deepen worthlessness
    \item \textbf{Social factors:} isolation, discrimination, and bullying attack a person's sense of belonging and value
    \item \textbf{Chronic illness and disability:} pain, fatigue, and loss of function can make people feel like a burden
    \item \textbf{Perfectionism:} unrelenting high standards and self-blame when they are not met
    \item \textbf{Cognitive patterns:} a negative bias that filters out achievements and magnifies failures
    \item \textbf{Substance use:} alcohol and drugs worsen mood and can be used to escape these feelings, creating a cycle
\end{itemize}

\section*{Clinical Features}
\begin{itemize}
    \item Persistent self-criticism and negative self-talk ("I'm useless", "I'm a burden")
    \item Low mood, loss of interest, and reduced energy
    \item Guilt, shame, and excessive blame of oneself for events
    \item Difficulty accepting praise or recognising achievements
    \item Withdrawal from others and avoidance of social contact
    \item Loss of motivation, poor concentration, and changes in sleep and appetite
    \item Hopelessness about the future and, importantly, thoughts of self-harm or suicide
    \item Physical symptoms such as fatigue, tension, and aches
\end{itemize}

\section*{Differential Diagnosis}
The feeling of worthlessness is usually part of a broader picture.
\begin{itemize}
    \item Major depressive disorder, the most common context
    \item Persistent depressive disorder (dysthymia) with long-standing low self-esteem
    \item Anxiety disorders and social anxiety
    \item Post-traumatic stress disorder after trauma or abuse
    \item Adjustment disorder after a major life stressor
    \item Bipolar disorder, in which worthlessness may occur in depressive phases
    \item Medical conditions affecting mood, including thyroid disease, and the psychological impact of chronic illness
    \item Normal, transient low mood in response to setbacks, which resolves with support
\end{itemize}

\section*{When to Seek Medical Care}
Anyone experiencing persistent worthlessness, especially with low mood, should speak to a doctor or mental health professional. Early help prevents the feelings from deepening and reduces the risk of crisis.

\textbf{Contact your local emergency services for:}
\begin{itemize}
    \item Thoughts of harming yourself or ending your life
    \item Making plans or taking steps to harm yourself
    \item Feeling that you cannot keep yourself safe
    \item Any immediate danger to yourself or others
\end{itemize}

\textbf{Support:} local crisis, youth, and mental-health services may be available around the clock. Contact local emergency services if there is immediate danger.

\section*{Diagnosis and Investigations}
\begin{itemize}
    \item \textbf{Clinical interview:} a doctor or psychologist explores the feelings, their onset, triggers, and impact, and asks directly about self-harm and suicide risk
    \item \textbf{Depression and anxiety screening tools:} structured questionnaires help assess severity and track progress
    \item \textbf{Assessment of alcohol and drug use:} substances may be driving or worsening the mood
    \item \textbf{Physical health review:} thyroid function and other tests if a medical cause is possible
    \item \textbf{Risk assessment:} establishing a safety plan whenever suicidal thoughts are present
    \item \textbf{Review of supports:} understanding the person's social and practical situation informs the care plan
\end{itemize}

\section*{Management}
\begin{itemize}
    \item \textbf{Psychological therapy:} cognitive behavioural therapy (CBT) directly challenges the distorted thoughts that maintain worthlessness; acceptance and commitment therapy (ACT) and interpersonal therapy are also effective
    \item \textbf{Antidepressant medicines:} helpful when worthlessness is part of moderate to severe depression, usually in combination with therapy
    \item \textbf{Safety planning:} a written plan identifying warning signs, coping strategies, and people to contact in a crisis
    \item \textbf{Behavioural activation:} gradually reintroducing meaningful activities and social contact, which rebuilds a sense of value
    \item \textbf{Treating underlying conditions:} anxiety, PTSD, substance use, and chronic pain all respond to specific treatment
    \item \textbf{Peer support:} support groups and online communities reduce isolation and shame
    \item \textbf{Family and practical support:} addressing financial, housing, and relationship stressors
\end{itemize}

\section*{Complications}
\begin{itemize}
    \item Self-harm and suicide, the most serious risks
    \item Social withdrawal and breakdown of relationships
    \item Impaired work or study performance and financial consequences
    \item Alcohol and drug misuse as an attempt to cope
    \item Worsening and prolongation of depression if untreated
    \item Physical health effects of chronic stress and poor self-care
\end{itemize}

\section*{Prognosis and Outcomes}
Feelings of worthlessness respond well to treatment, and the outlook is very good for most people. Therapy and, where needed, medication produce meaningful improvement in weeks to months, and the negative self-beliefs that sustain the feelings can be permanently reshaped. Relapses can occur, but people who learn to recognise early warning signs and maintain their support network manage them successfully. The most important message is that worthlessness is a symptom of illness and stress, not a true reflection of a person's value, and it passes with the right help.

\section*{Prevention and Self-Care}
\begin{itemize}
    \item Challenge negative self-talk by asking whether you would say the same words to a friend
    \item Keep connected: talk to trusted people, even when you feel like withdrawing
    \item Maintain routines of sleep, meals, and physical activity
    \item Set small, achievable goals and acknowledge each accomplishment
    \item Limit alcohol, which deepens low mood
    \item Seek professional help early rather than waiting for the feelings to pass
    \item If you are worried about someone, ask them directly how they are feeling and whether they have thought of harming themselves; asking is safe and opens the door to help
\end{itemize}

\section*{Sources and Further Reading}
The following reputable sources provide further information for healthcare students and patients seeking reliable, current advice about feeling worthless and low self-esteem.
\begin{itemize}
    \item Beyond Blue. \textit{Depression and low self-esteem.} https://www.beyondblue.org.au/
    \item Healthdirect Australia. \textit{Low self-esteem and depression.} https://www.healthdirect.gov.au/
    \item Lifeline. \textit{Crisis support and resources.} https://www.lifeline.org.au/
    \item Black Dog Institute. \textit{Depression resources.} https://www.blackdoginstitute.org.au/
    \item NHS. \textit{Low self-esteem.} https://www.nhs.uk/mental-health/
    \item Mayo Clinic. \textit{Depression.} https://www.mayoclinic.org/
\end{itemize}
